% Chalkline handout: Lecture Notes (preset lecture-notes v1).
% Document 01JCA000000000000000000017, version 1.
% Compile with LuaLaTeX, TeX Live 2025 or later: latexmk -lualatex main.tex
\DocumentMetadata{lang=es, pdfversion=2.0, pdfstandard=ua-2, tagging=on, tagging-setup={math/setup=mathml-SE}}
\documentclass[11pt]{article}
\usepackage[letterpaper,margin=1in,headheight=15pt]{geometry}
\usepackage{amsmath}
\usepackage{fontspec}
\usepackage{unicode-math}
\directlua{luaotfload.add_fallback("clfallback", {"Latin Modern Math:mode=node;", "file:FreeSerif.otf:mode=node;"})}
\setmainfont{Latin Modern Roman}[RawFeature={fallback=clfallback}]
\setmathfont{Latin Modern Math}
\hyphenpenalty=10000
\exhyphenpenalty=10000
\emergencystretch=3em
\usepackage{xcolor}
\definecolor{clBlue}{HTML}{1F4E9C}
\definecolor{clRed}{HTML}{B3261E}
\definecolor{clGreen}{HTML}{1E6B34}
\definecolor{clOrange}{HTML}{A65200}
\definecolor{clPurple}{HTML}{6B3FA0}
\definecolor{clBrown}{HTML}{7A4A1E}
\usepackage{graphicx}
\usepackage{array}
\usepackage{booktabs}
\usepackage{tcolorbox}
\usepackage{fancyhdr}
\usepackage{lastpage}
\pagestyle{fancy}
\fancyhf{}
\renewcommand{\headrulewidth}{0.4pt}
\NewDocumentCommand{\clheader}{mm}{\fancyhead[L]{\small #1}\fancyhead[R]{\small #2}}
\NewDocumentCommand{\clfooter}{m}{\fancyfoot[C]{\small #1}}
\NewDocumentCommand{\cltitle}{mm}{%
  \title{#1}\author{#2}\date{}%
  \maketitle\thispagestyle{fancy}\vspace{-2em}}
\NewDocumentCommand{\clquizheader}{mmm}{%
  \par\noindent #1: \rule{6cm}{0.4pt}\hfill #2: \rule{3cm}{0.4pt}\hfill #3: \rule{2cm}{0.4pt}\par\bigskip}
\NewDocumentEnvironment{cldefinition}{mm}{\begin{tcolorbox}[colback=white,colframe=black!55,boxrule=0.6pt,arc=2pt,left=6pt,right=6pt,top=4pt,bottom=4pt,before skip=8pt,after skip=8pt]\noindent\textbf{#1\IfBlankF{#2}{: #2}.}\ \ignorespaces}{\end{tcolorbox}}
\NewDocumentEnvironment{cltheorem}{mm}{\begin{tcolorbox}[colback=white,colframe=black!55,boxrule=0.6pt,arc=2pt,left=6pt,right=6pt,top=4pt,bottom=4pt,before skip=8pt,after skip=8pt]\noindent\textbf{#1\IfBlankF{#2}{: #2}.}\ \ignorespaces}{\end{tcolorbox}}
\NewDocumentEnvironment{clnote}{mm}{\par\medskip\noindent\textbf{#1\IfBlankF{#2}{: #2}.}\ \ignorespaces}{\par\medskip}
\NewDocumentEnvironment{clexample}{mm}{\par\medskip\noindent\textbf{#1\IfBlankF{#2}{: #2}.}\ \ignorespaces}{\par\medskip}
\NewDocumentEnvironment{clannouncement}{mm}{\par\medskip\noindent\textbf{#1\IfBlankF{#2}{: #2}.}\ \ignorespaces}{\par\medskip}
\NewDocumentEnvironment{clproblem}{mm}{\par\medskip\noindent\textbf{#1.}\IfBlankF{#2}{ \textit{(#2)}}\ \ignorespaces}{\par}
\NewDocumentEnvironment{clparts}{}{\begin{enumerate}[label=(\alph*)]}{\end{enumerate}}
\NewDocumentEnvironment{clsolution}{m}{\par\smallskip\noindent\textit{#1.}\ \ignorespaces}{\par\smallskip}
\NewDocumentCommand{\clanswer}{mm}{\par\smallskip\noindent\textit{#1:}\ #2\par}
\NewDocumentCommand{\clanswerboxed}{mm}{\par\smallskip\noindent\textit{#1:}\ \fbox{#2}\par}
\NewDocumentCommand{\clanswerspace}{m}{\par\vspace{#1\baselineskip}}
\NewDocumentCommand{\clplaceholder}{m}{\fbox{\textit{#1}}}
\setlength{\parindent}{0pt}
\setlength{\parskip}{0.5\baselineskip}
\usepackage{hyperref}
\hypersetup{hidelinks}
\hypersetup{pdftitle={Derivadas e integrales}, pdfdisplaydoctitle=true}
\clheader{Cálculo I}{21 de octubre de 2026}
\clfooter{Página \thepage\ de \pageref{LastPage}}

\begin{document}
\cltitle{Derivadas e integrales}{Cálculo I · Ana Müller-Šimić · 21 de octubre de 2026}

\section*{La derivada}

% blk:01JCA000000000000000000002
\[
\frac{dy}{dx}= \lim_{h \to 0}\frac{f(x + h) - f(x)}{h}
\]

% blk:01JCA000000000000000000004
\[
\frac{d}{dx}x^{n} = n x^{n-1}
\]

% blk:01JCA000000000000000000006
\[
f'(x) = 3x^{2} \cdot 2
\]

% blk:01JCA00000000000000000000A
\begin{cltheorem}{Teorema}{Regla de la cadena}
% blk:01JCA000000000000000000008
\[
\frac{d}{dx}f(g(x)) = f'(g(x)) \cdot g'(x)
\]
\end{cltheorem}

% blk:01JCA00000000000000000000C
La velocidad \(\vec{v}= \vec{v}_{0} + \vec{a}t\) con \(a = 9.81\) y \(\mathbf{F}= m \mathbf{a}\).

% blk:01JCA00000000000000000000E
\textcolor{clRed}{\textbf{Importante: }}no olvides la constante \textcolor{clBlue}{\(C\)}.

\section*{La integral}

% blk:01JCA00000000000000000000H
\[
\int_{0}^{1} x^{2} dx = \frac{1}{3}
\]

% blk:01JCA00000000000000000000K
\[
\int_{0}^{\pi}\sin x \, dx = 2
\]

% blk:01JCA00000000000000000000N
\[
A = 3.14 \times r^{2}
\]

% blk:01JCA00000000000000000000Q
\[
f(x) = \begin{cases}x^{2}&x \geq 0 \\ -x&x < 0\end{cases}
\]

% blk:01JCA00000000000000000000X
\begin{clproblem}{1}{}
% blk:01JCA00000000000000000000S
Calcula \(\int_{1}^{2} \frac{1}{x}dx\).

% blk:01JCA00000000000000000000V
\clanswerboxed{Respuesta}{\(\ln 2 \approx 0.693\)}
\end{clproblem}

% blk:01JCA000000000000000000015
\begin{clproblem}{2}{}
% blk:01JCA00000000000000000000Z
Resuelve el sistema:

% blk:01JCA000000000000000000011
\[
\begin{pmatrix}1 & 2 \\ 3 & 4\end{pmatrix} \begin{pmatrix}x \\ y\end{pmatrix} = \begin{pmatrix}5 \\ 6\end{pmatrix}
\]
% blk:01JCA000000000000000000013
\clanswerboxed{Respuesta}{\(x = -4\), \(y = 4.5\)}
\end{clproblem}

\end{document}
